\documentclass[12pt,brazil]{article}
\usepackage[utf8]{inputenc}
\usepackage[brazil]{babel}

\title{O gavião-real e a floresta}
\author{Equipe Harpia}

\begin{document}
\maketitle

\begin{abstract}
Este artigo de exemplo descreve a dieta do gavião-real e serve para ensinar a revisar textos
científicos com a biblioteca Harpia. Ele contém erros de propósito.
\end{abstract}

\section{Introdução}

O gavião-real (\textit{Harpia harpyja}) é uma das maiores aves de rapina do mundo
\cite{silva2020}. Ele vive nas florestas tropicais da América do Sul e caça preguiças e macacos.
Vale ressaltar que a espécie depende de árvores altas para fazer seus ninhos.

A conservação da espécie depende de de grandes áreas contínuas de floresta \cite{souza2018}.

% A frase abaixo é longa demais de propósito.
Nos estudos de campo realizados ao longo de vários anos em diferentes regiões da Amazônia, os
pesquisadores observaram que os casais de gaviões retornam ao mesmo ninho em temporadas
sucessivas, alimentam o filhote durante muitos meses e defendem o território contra outras aves
de rapina, o que torna a perda de uma única árvore um problema sério para toda a população local.

\section{Método}

A harpia é uma ave de rapina neotropical de grande porte, encontrada desde o sul do México até
o norte da Argentina, que se alimenta principalmente de mamíferos arborícolas como preguiças e
macacos capturados no dossel da floresta \cite{lima2015}.

Comparamos registros de ninhos com mapas de desmatamento.


\section{Conclusão}

Proteger o gavião-real significa proteger a floresta inteira \cite{inexistente2021}.

\bibliographystyle{plain}
\bibliography{refs}
\end{document}
